% Lean: EconLean/Macro/DP/Doeblin.lean, EconLean/Macro/DP/ParamInvariant.lean,
%       EconLean/Macro/HA/Aiyagari.lean, doeblin_pow in EconLean/Macro/DP/MarkovChain.lean
\chapter{Doeblin, Parametric Stationary Distributions and Aiyagari Equilibrium}
\label{chap:aiyagari}

Sources: Stokey--Lucas (SL) \S11.4, Thm~11.12, Thm~12.13; Meyn--Tweedie Thm~16.2.4;
Aiyagari (1994); Ljungqvist--Sargent ch.~18; Hopenhayn--Prescott (1992); A\c{c}{\i}kg\"oz (2018).

\section{Doeblin's theorem on a general state space}
\label{sec:doeblin_general}

Throughout, $K$ is a Markov kernel on a measurable space $X$. It satisfies the
\emph{minorization} with constant $\varepsilon\in[0,1]$ and probability measure $\varphi$ if
$K(x,A) \ge \varepsilon\,\varphi(A)$ for every $x$ and measurable $A$.

\begin{lemma}[Bounded functions are integrable]
\label{lem:integrable_of_bounded}
\lean{Econ.integrable_of_bounded}
\leanok
A measurable $f$ with $a \le f \le b$ is integrable against any finite measure.
\end{lemma}
\begin{proof}
\leanok
$|f| \le \max(|a|,|b|)$; apply \texttt{Integrable.of\_bound}.
\end{proof}

\begin{lemma}[Integral bound from a Hahn decomposition]
\label{lem:integral_sub_le_hahn}
\lean{Econ.integral_sub_le_hahn}
\leanok
Let $\mu,\nu$ be finite and $P$ measurable with $\nu \le \mu$ on subsets of $P$ and $\mu\le\nu$ on subsets
of $P^c$. If $f$ is measurable with $0 \le f \le c$, then $\int f\,d\mu - \int f\,d\nu \le c\,(\mu(P) - \nu(P))$.
\end{lemma}
\begin{proof}
\leanok
\uses{lem:integrable_of_bounded}
Split both integrals over $P$ and $P^c$ (\texttt{integral\_add\_compl}). On $P^c$, $\mu|_{P^c} \le \nu|_{P^c}$
and $f\ge0$ give $\int_{P^c} f\,d\mu \le \int_{P^c} f\,d\nu$ (\texttt{integral\_mono\_measure}). On $P$, write
$\mu|_P = \rho + \nu|_P$ with $\rho = \mu|_P - \nu|_P \ge 0$ (\texttt{Measure.sub\_add\_cancel\_of\_le}). Then
$\int_P f\,d\mu - \int_P f\,d\nu = \int f\,d\rho \le c\,\rho(X) = c(\mu(P) - \nu(P))$.
\end{proof}

\begin{lemma}[Composition as an integral]
\label{lem:comp_real_eq_integral}
\lean{Econ.comp_real_eq_integral}
\leanok
$(K\circ\rho)(A) = \int K(x,A)\,d\rho(x)$ for a probability $\rho$ and measurable $A$.
\end{lemma}
\begin{proof}
\leanok
Apply \texttt{integral\_comp\_measure} (Lemma~\ref{lem:feller_integrals}) to the indicator of $A$.
\end{proof}

\begin{lemma}[One Doeblin step]
\label{lem:doeblin_step}
\lean{Econ.doeblin_step}
\leanok
Under the minorization, if $\mu(A) - \nu(A) \le c$ for all measurable $A$ (probabilities $\mu,\nu$), then
$(K\circ\mu)(A) - (K\circ\nu)(A) \le (1-\varepsilon)c$ for all measurable $A$.
\end{lemma}
\begin{proof}
\leanok
\uses{lem:integral_sub_le_hahn, lem:comp_real_eq_integral}
Take a Hahn decomposition $P$ of $\mu - \nu$ (\texttt{hahn\_decomposition}). Fix $A$ and let
$g(x) = K(x,A) - \varepsilon\varphi(A)$. Minorization at $A$ gives $g\ge0$. Minorization at $A^c$ gives
$1 - K(x,A) \ge \varepsilon(1-\varphi(A))$, so $g \le 1 - \varepsilon$. By Lemma~\ref{lem:comp_real_eq_integral},
$\int g\,d\rho = (K\circ\rho)(A) - \varepsilon\varphi(A)$, so the constant cancels in the difference. By
Lemma~\ref{lem:integral_sub_le_hahn}, the difference is at most $(1-\varepsilon)(\mu(P)-\nu(P)) \le (1-\varepsilon)c$.
\end{proof}

\begin{definition}[Iterated distribution]
\label{def:iterComp}
\lean{Econ.iterComp, Econ.isProbabilityMeasure_iterComp}
\leanok
$\mu K^0 = \mu$, $\mu K^{n+1} = K \circ (\mu K^n)$; each is a probability measure.
\end{definition}

\begin{theorem}[Doeblin: geometric convergence; SL Thm 11.12, minorization case]
\label{thm:doeblin_convergence}
\lean{Econ.doeblin_convergence}
\leanok
\uses{def:iterComp}
Under the minorization, if $\pi$ is $K$-invariant, then for every probability $\mu$, every $n$ and
measurable $A$, $|\mu K^n(A) - \pi(A)| \le (1-\varepsilon)^n$.
\end{theorem}
\begin{proof}
\leanok
\uses{lem:doeblin_step}
Induct on $n$, carrying both one-sided bounds $\mu K^n(A) - \pi(A) \le (1-\varepsilon)^n$ and
$\pi(A) - \mu K^n(A) \le (1-\varepsilon)^n$. At $n=0$ both are differences of numbers in $[0,1]$. The step
is Lemma~\ref{lem:doeblin_step} applied to $(\mu K^n, \pi)$ and to $(\pi, \mu K^n)$, using $K\circ\pi = \pi$.
\end{proof}

\begin{theorem}[Doeblin: uniqueness]
\label{thm:doeblin_unique}
\lean{Econ.doeblin_unique}
\leanok
Under the minorization with $\varepsilon > 0$, any two invariant probabilities coincide.
\end{theorem}
\begin{proof}
\leanok
\uses{thm:doeblin_convergence}
Starting from $\pi'$, $\pi' K^n = \pi'$, so $|\pi'(A) - \pi(A)| \le (1-\varepsilon)^n \to 0$. Equal values on
measurable sets give equal measures.
\end{proof}

\begin{lemma}[Kernel powers]
\label{lem:kernel_pow}
\lean{Econ.isMarkovKernel_pow, Econ.invariant_pow}
\leanok
$K^n$ is Markov, and a $K$-invariant measure is $K^n$-invariant.
\end{lemma}
\begin{proof}
\leanok
Induction with $K^{n+1} = K^n \circ K$; Mathlib's \texttt{Kernel.Invariant.comp} and
\texttt{Measure.id\_comp}.
\end{proof}

\begin{theorem}[Uniqueness under an $N$-step minorization; SL Condition D]
\label{thm:unique_invariant_of_pow}
\lean{Econ.unique_invariant_of_pow}
\leanok
If $K^N(x,\cdot) \ge \varepsilon\varphi$ for all $x$, with $\varepsilon>0$, then $K$ has at most one invariant
probability.
\end{theorem}
\begin{proof}
\leanok
\uses{lem:kernel_pow, thm:doeblin_unique}
Both invariant measures are $K^N$-invariant (Lemma~\ref{lem:kernel_pow}); apply
Theorem~\ref{thm:doeblin_unique} to $K^N$.
\end{proof}

\begin{lemma}[Iterates along powers]
\label{lem:iterComp_pow}
\lean{Econ.iterComp_add_pow, Econ.iterComp_pow}
\leanok
\uses{def:iterComp}
$\mu K^{m+N} = K^N\circ(\mu K^m)$, and $\mu K^{a + Nb} = (\mu K^a)(K^N)^b$.
\end{lemma}
\begin{proof}
\leanok
Induction on $N$ using $K^{N+1} = K\circ K^N$ and \texttt{Measure.comp\_assoc}; then induction on $b$.
\end{proof}

\begin{theorem}[Geometric convergence under an $N$-step minorization]
\label{thm:doeblin_convergence_pow}
\lean{Econ.doeblin_convergence_pow}
\leanok
If $K^N(x,\cdot)\ge\varepsilon\varphi$ for all $x$ and $\pi$ is $K$-invariant, then for every probability $\mu$,
$|\mu K^n(A) - \pi(A)| \le (1-\varepsilon)^{\lfloor n/N\rfloor}$.
\end{theorem}
\begin{proof}
\leanok
\uses{lem:iterComp_pow, lem:kernel_pow, thm:doeblin_convergence}
Write $n = (n \bmod N) + N\lfloor n/N\rfloor$, so $\mu K^n = (\mu K^{n\bmod N})(K^N)^{\lfloor n/N\rfloor}$
(Lemma~\ref{lem:iterComp_pow}). $\pi$ is $K^N$-invariant (Lemma~\ref{lem:kernel_pow}); apply
Theorem~\ref{thm:doeblin_convergence} to $K^N$.
\end{proof}

\begin{lemma}[Powers of stochastic matrices]
\label{lem:stochastic_pow}
\lean{Econ.IsStochastic.pow, Econ.IsInvariant.pow}
\leanok
\uses{def:markov}
$Q^n$ is stochastic, and a $Q$-invariant distribution is $Q^n$-invariant.
\end{lemma}
\begin{proof}
\leanok
Induction on $n$. Entries of $Q^{n}Q$ are sums of products of nonnegative numbers. Row sums:
exchange the sums and use that rows of $Q$ and $Q^n$ sum to one.
\end{proof}

\begin{theorem}[Finite Doeblin, $N$-step form]
\label{thm:doeblin_pow}
\lean{Econ.doeblin_pow}
\leanok
\uses{def:markov}
If $Q$ is stochastic and every entry of $Q^N$ is at least $\varepsilon>0$, then $Q$ has a unique invariant
distribution $\pi$ and $\|\mu Q^n - \pi\|_1 \le 2(1 - |Z|\varepsilon)^{\lfloor n/N\rfloor}$. Example: $Q =
\begin{pmatrix}0&1\\1/2&1/2\end{pmatrix}$ has a zero entry but $Q^2 \ge 1/4$ entrywise.
\end{theorem}
\begin{proof}
\leanok
\uses{lem:stochastic_pow, thm:doeblin, thm:exists_invariant}
Theorem~\ref{thm:doeblin} applies to $Q^N$ (stochastic by Lemma~\ref{lem:stochastic_pow}) and gives a unique
$Q^N$-invariant $\pi_0$. A $Q$-invariant $\pi$ exists (Theorem~\ref{thm:exists_invariant}) and is
$Q^N$-invariant, so $\pi = \pi_0$; the same argument gives uniqueness. For convergence write
$Q^n = Q^{n \bmod N}(Q^N)^{\lfloor n/N\rfloor}$ and apply the $Q^N$ bound to the distribution $\mu Q^{n\bmod N}$.
\end{proof}

\section{Continuity of invariant distributions in a parameter}
\label{sec:param_invariant}

$X$ is a compact metrizable space with its Borel $\sigma$-algebra; $\Theta$ is a topological space;
probability measures carry the weak topology.

\begin{definition}[Jointly Feller family]
\label{def:jointly_feller}
\lean{Econ.IsJointlyFeller}
\leanok
\uses{def:feller}
$(\kappa_\theta)$ is jointly Feller if $(\theta,x) \mapsto \int f\,d\kappa_\theta(x)$ is jointly continuous for every
bounded continuous $f$.
\end{definition}

\begin{lemma}[Limits of invariant distributions are invariant]
\label{lem:invariant_of_tendsto}
\lean{Econ.invariant_of_tendsto, Econ.isClosed_invariant_graph}
\leanok
\uses{def:jointly_feller}
If, along a nontrivial filter, $\theta_i\to\theta_0$, $\mu_i\to\mu_0$ weakly and eventually $\mu_i$ is
$\kappa_{\theta_i}$-invariant, then $\mu_0$ is $\kappa_{\theta_0}$-invariant. Hence the graph
$\{(\theta,\mu) : \mu \text{ is } \kappa_\theta\text{-invariant}\}$ is closed.
\end{lemma}
\begin{proof}
\leanok
\uses{lem:feller_integrals}
Fix bounded continuous $f$ and let $G_\theta(x) = \int f\,d\kappa_\theta(x)$. Currying the jointly continuous $G$
gives a continuous map into $C(X)$; on compact $X$ this is uniform convergence $G_\theta\to G_{\theta_0}$
(\texttt{ContinuousMap.tendsto\_iff\_tendstoUniformly}). Eventually invariance gives
$\int f\,d\mu_i = \int G_{\theta_i}\,d\mu_i$. The error $\int (G_{\theta_i} - G_{\theta_0})\,d\mu_i$ is at most the
uniform distance, so tends to $0$; and $\int G_{\theta_0}\,d\mu_i \to \int G_{\theta_0}\,d\mu_0$,
$\int f\,d\mu_i\to\int f\,d\mu_0$. Hence $\int f\,d\mu_0 = \int G_{\theta_0}\,d\mu_0 = \int f\,d(\kappa_{\theta_0}\circ\mu_0)$,
and bounded continuous functions determine finite measures. For the graph, use
\texttt{isClosed\_iff\_ultrafilter} with the two projections.
\end{proof}

\begin{theorem}[SL Thm 12.13]
\label{thm:tendsto_invariant_of_unique}
\lean{Econ.tendsto_invariant_of_unique}
\leanok
\uses{def:jointly_feller}
Let $(\kappa_\theta)$ be jointly Feller Markov kernels and $\pi_\theta$ any $\kappa_\theta$-invariant probabilities.
If $\kappa_{\theta_0}$ has a unique invariant probability, then $\pi_\theta \to \pi_{\theta_0}$ weakly as
$\theta\to\theta_0$.
\end{theorem}
\begin{proof}
\leanok
\uses{lem:invariant_of_tendsto}
The space of probability measures on $X$ is compact (Prokhorov), so by
\texttt{tendsto\_nhds\_of\_unique\_mapClusterPt} it suffices to show every cluster point $\mu$ equals
$\pi_{\theta_0}$. Take an ultrafilter $U\le\mathcal N(\theta_0)$ along which $\pi_\theta\to\mu$; by
Lemma~\ref{lem:invariant_of_tendsto} $\mu$ is $\kappa_{\theta_0}$-invariant, and uniqueness gives $\mu=\pi_{\theta_0}$.
\end{proof}

\begin{corollary}[Continuity of the invariant distribution and its integrals]
\label{cor:continuous_invariant}
\lean{Econ.continuous_invariant_of_unique, Econ.continuous_integral_invariant}
\leanok
If every $\kappa_\theta$ has a unique invariant $\pi_\theta$, then $\theta\mapsto\pi_\theta$ is continuous, and so is
$\theta \mapsto \int f\,d\pi_\theta$ for bounded continuous $f$.
\end{corollary}
\begin{proof}
\leanok
\uses{thm:tendsto_invariant_of_unique}
Pointwise continuity from Theorem~\ref{thm:tendsto_invariant_of_unique}; integration against bounded
continuous $f$ is weakly continuous.
\end{proof}

\section{Savings problems parametrized by prices}
\label{sec:param_savings}

\begin{definition}[Parametrized savings problem]
\label{def:param_savings}
\lean{Econ.ParamSavings, Econ.ParamSavings.slice, Econ.ParamSavings.S, Econ.ParamSavings.continuous_hS}
\leanok
\uses{def:stoch_savings}
As in Definition~\ref{def:stoch_savings}, but cash on hand $h(\theta, a, z)$ depends on a price $\theta$,
jointly continuously in $(\theta,a)$, with $h(\theta,\cdot,z)$ concave and nonnegative on $[0,\infty)$. The slice
at $\theta$ is the ordinary savings problem with $h(\theta,\cdot,\cdot)$. Aiyagari: $h(r,a,z) = (1+r)a + w(r)e(z)$.
\end{definition}

\begin{definition}[Augmented problem]
\label{def:param_savings_aug}
\lean{Econ.ParamSavings.augΓ, Econ.ParamSavings.augF, Econ.ParamSavings.augΓ_isClosed_graph,
  Econ.ParamSavings.augΓ_lhc, Econ.ParamSavings.aug}
\leanok
\uses{def:param_savings, def:stoch_dp, def:contCorr}
For compact Hausdorff $\Theta$, the stochastic DP on states $((\theta,a),z)$ with feasible set
$\{(\theta,a') : 0\le a'\le\min(\bar k, h(\theta,a,z))\}$ and return $u(h(\theta,a,z) - a')$. The price never changes.
The feasible correspondence has a closed graph (so is UHC into the compact $\Theta\times[0,\bar k]$) and is LHC
via the continuous selection $(\theta', a', z') \mapsto (\theta', \min(a', h))$; returns are bounded on the compact
state space.
\end{definition}

\begin{theorem}[The augmented problem restricts to its slices]
\label{thm:value_slice}
\lean{Econ.ParamSavings.value_slice, Econ.ParamSavings.objective_slice, Econ.ParamSavings.policy_slice}
\leanok
\uses{def:param_savings_aug, thm:stoch_bellman}
$V^{\mathrm{aug}}((\theta,a),z) = V_\theta(a,z)$, and the augmented policy at $((\theta,a),z)$ is
$\{\theta\}\times G_\theta(a,z)$.
\end{theorem}
\begin{proof}
\leanok
The function $w(a,z) = V^{\mathrm{aug}}((\theta,a),z)$ is bounded continuous. The augmented Bellman equation at
$((\theta,a),z)$ maximizes over the image of the slice's feasible set under $a'\mapsto(\theta,a')$, so $w$ solves
the slice's Bellman equation, and $w = V_\theta$ by uniqueness (\texttt{value\_unique}). Then the augmented
objective at $(\theta,a')$ equals the slice objective at $a'$, and the argmax sets correspond.
\end{proof}

\begin{theorem}[Savings policies are jointly continuous in prices]
\label{thm:exists_joint_policy}
\lean{Econ.ParamSavings.exists_joint_policy}
\leanok
\uses{def:param_savings}
There is $g$ with $G_\theta(s) = \{g(\theta,s)\}$ for all $\theta, s$, and $(\theta,a,z)\mapsto g(\theta,(a,z))$
jointly continuous.
\end{theorem}
\begin{proof}
\leanok
\uses{thm:stoch_value_concave, thm:value_slice, cor:continuous_argmax_singleton}
Each slice has a unique policy $g_\theta$ (concavity, Theorem~\ref{thm:stoch_value_concave}). By
Theorem~\ref{thm:value_slice} the augmented argmax is the singleton $\{(\theta, g_\theta(a,z))\}$. Berge
(Corollary~\ref{cor:continuous_argmax_singleton}) makes it continuous; take the second component.
\end{proof}

\begin{lemma}[The policy kernels are jointly Feller]
\label{lem:isJointlyFeller_policyKernel}
\lean{Econ.ParamSavings.isJointlyFeller_policyKernel}
\leanok
\uses{def:jointly_feller, def:policy_kernel}
If $g$ is jointly continuous, the kernels $\kappa_\theta((a,z),\cdot) = \sum_{z'} Q(z,z')\delta_{(g(\theta,a,z),z')}$
are jointly Feller.
\end{lemma}
\begin{proof}
\leanok
\uses{lem:policy_kernel}
$\int f\,d\kappa_\theta(a,z) = \sum_{z'}Q(z,z') f(g(\theta,a,z),z')$, a finite sum of continuous functions.
\end{proof}

\begin{definition}[Assets and the uniqueness hypothesis]
\label{def:unique_stationary}
\lean{Econ.ParamSavings.assets, Econ.ParamSavings.UniqueStationary, Econ.ParamSavings.assets_le,
  Econ.ParamSavings.integral_assets_le, Econ.ParamSavings.exists_stationary_of_policy}
\leanok
\uses{def:param_savings, def:policy_kernel}
$\mathrm{assets}(a,z) = a \in [0,\bar k]$, so $\int \mathrm{assets}\,d\mu \le \bar k$. \emph{UniqueStationary}$(\theta)$:
the optimal policy's chain at price $\theta$ has at most one invariant probability. At least one always
exists (\texttt{exists\_stationary\_of\_policy}: the policy equals the continuous one of
Theorem~\ref{thm:stoch_savings_stationary}, whose kernel is Feller). This holds, e.g., under an
$N$-step Doeblin minorization (Theorem~\ref{thm:unique_invariant_of_pow}) or Hopenhayn--Prescott monotone
mixing (not formalized).
\end{definition}

\begin{theorem}[Aggregate asset supply is continuous in prices]
\label{thm:continuous_aggregate}
\lean{Econ.ParamSavings.continuous_aggregate}
\leanok
\uses{def:unique_stationary}
If UniqueStationary holds at every price, there are a jointly continuous policy $g$, invariant $\pi_\theta$, with
$\theta\mapsto\pi_\theta$ continuous and $K_s(\theta) = \int a\,d\pi_\theta$ continuous.
\end{theorem}
\begin{proof}
\leanok
\uses{thm:exists_joint_policy, lem:isJointlyFeller_policyKernel, thm:exists_invariant_of_feller,
  cor:continuous_invariant}
Theorem~\ref{thm:exists_joint_policy} gives $g$; each $g_\theta$ is continuous, so each kernel is Feller and has
an invariant $\pi_\theta$ (Theorem~\ref{thm:exists_invariant_of_feller}). The family is jointly Feller
(Lemma~\ref{lem:isJointlyFeller_policyKernel}); apply Corollary~\ref{cor:continuous_invariant}.
\end{proof}

\begin{theorem}[Asset-market clearing in Aiyagari's model; existence without uniqueness]
\label{thm:aiyagari_equilibrium}
\lean{Econ.ParamSavings.exists_asset_market_clearing}
\leanok
\uses{def:unique_stationary}
Prices range over $[r_{lo}, r_{hi}]$. Assume $K_d$ continuous, $K_d(r_{lo})\ge\bar k$, and $K_d(r_{hi}) \le \int a\,d\mu$
for some stationary $\mu$ at $r_{hi}$. Then some $r$ and stationary $\pi$ at $r$ satisfy $\int a\,d\pi = K_d(r)$.
Gaps: $K_d$ is an arbitrary demand curve (no firm, wage or labor-market clearing); the upper boundary condition
and the asset cap $\bar k$ are assumed, not derived; if the cap binds, the equilibrium need not be one of the
uncapped economy.
\end{theorem}
\begin{proof}
\leanok
\uses{thm:exists_joint_policy, lem:isJointlyFeller_policyKernel, lem:invariant_of_tendsto,
  thm:exists_invariant_of_feller}
Fix the jointly continuous policy (Theorem~\ref{thm:exists_joint_policy}); the kernels are jointly Feller and each
has an invariant probability. Let $C$ be the (closed, Lemma~\ref{lem:invariant_of_tendsto}) graph of invariant
probabilities, and $A$ (resp.\ $B$) the projection to prices of $C \cap \{\int a\,d\mu \le K_d(r)\}$ (resp.\ $\ge$).
Projection along the compact space of probability measures is closed (\texttt{isClosedMap\_fst\_of\_compactSpace}),
so $A,B$ are closed; they cover the interval, $r_{lo}\in A$ by \texttt{integral\_assets\_le}, $r_{hi}\in B$ by
hypothesis. Connectedness (\texttt{isPreconnected\_closed\_iff}) gives $r\in A\cap B$: stationary $\mu_1,\mu_2$ with
$k_1 \le K_d(r) \le k_2$. If $k_1 = k_2$ take $\mu_1$. Otherwise $t = (k_2 - K_d(r))/(k_2-k_1)\in[0,1]$ and
$t\mu_1 + (1-t)\mu_2$ is a stationary probability (\texttt{Measure.comp\_add}, \texttt{Measure.comp\_smul}) with
$\int a = tk_1 + (1-t)k_2 = K_d(r)$.
\end{proof}

\begin{remark}[Non-vacuity]
\label{rem:aiyagari_examples}
\lean{Econ.ParamSavings.aiyagariFamily, Econ.ParamSavings.noReturn, Econ.ParamSavings.noReturn_value,
  Econ.ParamSavings.noReturn_policy, Econ.ParamSavings.noReturn_uniqueStationary}
\leanok
\uses{thm:exists_joint_policy, thm:aiyagari_equilibrium, thm:doeblin_unique}
\emph{aiyagariFamily}: $u=\sqrt c$, $r\in[0,1/20]$, $h = (1+r)a + y(z)$, $y\in\{1/4,1/2\}$ i.i.d., $\beta = 9/10$;
Theorem~\ref{thm:exists_joint_policy} applies, and with $K_d(r) = 1 - 20r$ so does
Theorem~\ref{thm:aiyagari_equilibrium} (cheap boundary conditions $K_d(0)=\bar k$, $K_d(1/20) = 0$). \emph{noReturn}: $h = y(z)$ does not depend on assets. Both Bellman
equations at $(a,z)$ and $(0,z)$ maximize the same set, so $V$ does not depend on $a$; then $a'=0$ is optimal since
$\sqrt{\cdot}$ is monotone, and uniqueness forces $g = 0$. The policy kernel is the same measure from every state,
so the minorization holds with $\varepsilon=1$ and Theorem~\ref{thm:doeblin_unique} gives UniqueStationary
(degenerate in assets; non-degenerate uniqueness is open).
\end{remark}

\section{Aiyagari stationary general equilibrium}
\label{sec:aiyagari_ge}

\begin{definition}[Cobb--Douglas firm]
\label{def:cd_firm}
\lean{Econ.CDFirm, Econ.CDFirm.kOfR, Econ.CDFirm.wage, Econ.CDFirm.capDemand, Econ.CDFirm.output,
  Econ.CDFirm.one_sub_α_pos, Econ.CDFirm.kOfR_pos, Econ.CDFirm.wage_pos}
\leanok
$Y = K^\alpha L^{1-\alpha}$, $0<\alpha<1$, depreciation $\delta\ge0$. For $r > -\delta$:
$k(r) = (\alpha/(r+\delta))^{1/(1-\alpha)} > 0$, $w(r) = (1-\alpha)k(r)^\alpha > 0$, $K_d(r) = L\,k(r)$.
\end{definition}

\begin{lemma}[Firm optimality and zero profit]
\label{lem:cd_firm_foc}
\lean{Econ.CDFirm.foc_capital, Econ.CDFirm.foc_labor, Econ.CDFirm.output_eq, Econ.CDFirm.zero_profit}
\leanok
\uses{def:cd_firm}
$\alpha k(r)^{\alpha-1} = r+\delta$; $(1-\alpha)k(r)^\alpha = w(r)$; and for $L>0$,
$(r+\delta)K_d + wL = K_d^\alpha L^{1-\alpha}$.
\end{lemma}
\begin{proof}
\leanok
$k^{\alpha-1} = (\alpha/(r+\delta))^{\frac{\alpha-1}{1-\alpha}} = (\alpha/(r+\delta))^{-1}$ (\texttt{Real.rpow\_mul}).
Output at $K = Lk$ is $Lk^\alpha$ (\texttt{Real.mul\_rpow}, \texttt{Real.rpow\_add}). Then
$(r+\delta)Lk = \alpha k^{\alpha-1}\cdot kL = \alpha Lk^\alpha$ (\texttt{Real.rpow\_sub\_one}) and
$wL = (1-\alpha)Lk^\alpha$; they sum to $Lk^\alpha$.
\end{proof}

\begin{lemma}[Capital demand: continuity, monotonicity, divergence]
\label{lem:cd_capdemand}
\lean{Econ.CDFirm.continuousOn_kOfR, Econ.CDFirm.continuousOn_wage, Econ.CDFirm.continuousOn_capDemand,
  Econ.CDFirm.strictAntiOn_kOfR, Econ.CDFirm.tendsto_kOfR, Econ.CDFirm.exists_capDemand_ge}
\leanok
\uses{def:cd_firm}
On $r>-\delta$: $k$, $w$, $K_d$ are continuous; $k$ is strictly decreasing; $k(r)\to\infty$ as $r\downarrow-\delta$.
Hence for $L>0$, any $B$ and $r_0 > -\delta$ there is $r\in(-\delta, r_0]$ with $K_d(r)\ge B$.
\end{lemma}
\begin{proof}
\leanok
Continuity: \texttt{ContinuousOn.rpow\_const} with a nonnegative exponent. Monotonicity:
\texttt{Real.rpow\_lt\_rpow} with positive exponent. Divergence: $r+\delta\to0^+$, so
$\alpha/(r+\delta)\to\infty$ (\texttt{tendsto\_inv\_nhdsGT\_zero}), then \texttt{tendsto\_rpow\_atTop}.
The last claim intersects the eventual set $\{k \ge B/L\}$ with $(-\delta, r_0]$.
\end{proof}

\begin{lemma}[Income marginals and stationary assets]
\label{lem:aiyagari_marg}
\lean{Econ.ParamSavings.zfun, Econ.ParamSavings.zfun_apply, Econ.ParamSavings.marg,
  Econ.ParamSavings.integral_zfun, Econ.ParamSavings.marg_isInvariant, Econ.ParamSavings.integral_policy_eq}
\leanok
\uses{def:param_savings, def:policy_kernel, def:markov}
For any policy $g$ and any stationary $\mu$ of its kernel: the income marginal $q(z) = \mu(Z=z)$ is a
$Q$-invariant distribution, $\int F(z)\,d\mu = \sum_z F(z)q(z)$, and $\int g\,d\mu = \int a\,d\mu$.
\end{lemma}
\begin{proof}
\leanok
\uses{lem:policy_kernel, lem:feller_integrals}
Write $F(z) = \sum_z F(z)\mathbf 1_{\{z\}}$ and integrate (\texttt{integral\_finsetSum}). For invariance,
$q(z') = \int \mathbf 1_{z'}\,d(\kappa\circ\mu) = \int \sum_{z''}Q(z,z'')\mathbf 1_{z'}(z'')\,d\mu = \int Q(z,z')\,d\mu =
\sum_z q(z)Q(z,z')$. For assets, $\int a\,d\mu = \int a\,d(\kappa\circ\mu) = \int \sum_{z'}Q(z,z')g(a,z)\,d\mu =
\int g\,d\mu$ since rows of $Q$ sum to one.
\end{proof}

\begin{theorem}[A lower bound on savings by a one-period deviation]
\label{thm:le_policy}
\lean{Econ.ParamSavings.le_policy}
\leanok
\uses{def:param_savings}
If saving $m$ is always feasible, cash on hand is monotone in assets, and for every state, every $y<m$ and every
continuation choice $x(z')\le h(y,z')$,
$u(h-y) - u(h-m) < \beta\sum_{z'}Q(z,z')[u(h(m,z')-x(z')) - u(h(y,z')-x(z'))]$, then the optimal policy saves at
least $m$ everywhere.
\end{theorem}
\begin{proof}
\leanok
\uses{thm:stoch_bellman}
Suppose $y = g(a,z) < m$. Let $x(z') = g(y,z')$. Bellman: $V(a,z) = F(y) + \beta\sum Q V(y,\cdot)$ and
$V(y,z') = F(y,z',x) + \beta\sum QV(x,\cdot)$ (\texttt{value\_eq\_of\_mem\_policy}). Feasibility of $m$ today and of $x(z')$
from $m$ tomorrow (monotone $h$) gives $V(a,z) \ge F(m) + \beta\sum QV(m,\cdot)$ and
$V(m,z') \ge F(m,z',x) + \beta\sum QV(x,\cdot)$ (\texttt{value\_ge}). Subtracting yields
$F(y) - F(m) \ge \beta\sum Q[F(m,z',x) - F(y,z',x)]$, contradicting the hypothesis.
\end{proof}

\begin{definition}[Aiyagari economy]
\label{def:aiyagari_economy}
\lean{Econ.AiyagariEconomy, Econ.AiyagariEconomy.L, Econ.AiyagariEconomy.mem_Ioi,
  Econ.AiyagariEconomy.household, Econ.AiyagariEconomy.household_h, Econ.AiyagariEconomy.household_h_mono}
\leanok
\uses{def:param_savings, def:cd_firm, def:markov}
Households ($u,\beta,Q$, efficiency units $e\ge0$, cap $\bar k$), a Cobb--Douglas firm, prices
$r\in[r_{lo},r_{hi}]$ with $r_{lo} > -\delta$, $r_{lo}\ge-1$; $Q$ has a unique invariant $\bar q$ and
$L = \sum_z\bar q(z)e(z) > 0$. The household family has cash on hand $(1+r)a + w(r)e(z)$, monotone in $a$.
\end{definition}

\begin{theorem}[Existence of a stationary equilibrium with an exogenous asset cap]
\label{thm:aiyagari_ge}
\lean{Econ.AiyagariEconomy.exists_stationary_equilibrium}
\leanok
\uses{def:aiyagari_economy}
If $\bar k \le K_d(r_{lo})$ and at $r_{hi}$ some stationary distribution supplies at least $K_d(r_{hi})$, there are
$r$, the optimal policy $g$ at $(r,w(r))$, and a stationary $\pi$ with: $K = \int a\,d\pi = L\,k(r)$;
$\int e\,d\pi = L$; $\alpha(K/L)^{\alpha-1} = r+\delta$ and $w(r) = (1-\alpha)(K/L)^\alpha$; and
$\int c\,d\pi + \delta K = K^\alpha L^{1-\alpha}$ with $c = (1+r)a + we - g$. Household optimality is subject to
$a'\le\bar k$, and the cap need not be slack; deriving an endogenous cap from $\beta(1+r)<1$ is open.
\end{theorem}
\begin{proof}
\leanok
\uses{thm:aiyagari_equilibrium, lem:aiyagari_marg, lem:cd_firm_foc, lem:cd_capdemand}
Apply Theorem~\ref{thm:aiyagari_equilibrium} with $K_d = L\,k(\cdot)$ (continuous on the interval). Labor: by
Lemma~\ref{lem:aiyagari_marg} the income marginal of $\pi$ is $Q$-invariant, hence equals $\bar q$, so
$\int e\,d\pi = L$. Firm: $K/L = k(r)$ and Lemma~\ref{lem:cd_firm_foc}. Goods: $\int c = (1+r)K + wL - \int g =
rK + wL$ (Lemma~\ref{lem:aiyagari_marg}), and zero profit gives $rK + wL + \delta K = Y$.
\end{proof}

\begin{lemma}[Upper boundary condition from a savings bound]
\label{lem:hhi_of_le_policy}
\lean{Econ.AiyagariEconomy.hhi_of_le_policy}
\leanok
\uses{def:aiyagari_economy}
If at $r_{hi}$ every optimal policy saves at least $m \ge K_d(r_{hi})$ in every state, then some stationary
distribution at $r_{hi}$ supplies at least $K_d(r_{hi})$.
\end{lemma}
\begin{proof}
\leanok
\uses{lem:aiyagari_marg, def:unique_stationary}
Take a stationary $\mu$ (\texttt{exists\_stationary\_of\_policy}); $\int a\,d\mu = \int g\,d\mu \ge m$.
\end{proof}

\begin{theorem}[A capped economy as a non-vacuity witness]
\label{thm:sqrt_equilibrium}
\lean{Econ.sqrt_sub_le_div, Econ.div_le_sqrt_sub, Econ.CDFirm.half, Econ.CDFirm.half_kOfR,
  Econ.CDFirm.half_wage, Econ.half_isInvariant, Econ.AiyagariEconomy.sqrtEconomy,
  Econ.AiyagariEconomy.sqrt_L, Econ.AiyagariEconomy.sqrt_rhi, Econ.AiyagariEconomy.sqrt_saving,
  Econ.AiyagariEconomy.sqrt_equilibrium}
\leanok
\uses{thm:aiyagari_ge, thm:le_policy, lem:hhi_of_le_policy}
With $u=\sqrt c$, $\beta = 9/10$, i.i.d.\ two-state $Q$ with $e\equiv1$ ($L=1$), $\bar k = 1$, $\alpha=1/2$,
$\delta = 1/10$, $r\in[0,19/10]$: a stationary equilibrium of the capped economy exists. This only shows the
hypotheses of Theorem~\ref{thm:aiyagari_ge} are jointly satisfiable: there is no income risk, and since supply is at
most $\bar k=1$ every equilibrium has $r\ge 2/5$, so $\beta(1+r)>1$ and the cap binds. It is not a textbook Aiyagari
equilibrium, which has $r<1/\beta-1$.
\end{theorem}
\begin{proof}
\leanok
\uses{thm:doeblin}
$\bar q = (1/2,1/2)$ is the unique invariant distribution (Theorem~\ref{thm:doeblin}). $k(0) = 25 \ge \bar k$.
At $r_{hi} = 19/10$: $k = 1/16$, $w = 1/8$, $h = \tfrac{29}{10}a + \tfrac18$. Theorem~\ref{thm:le_policy} with
$m = 1/16$: with $d = m-y$, the cost today is $\sqrt{h-y}-\sqrt{h-m} \le d/(2\sqrt{h-m}) \le 2d$ since $h-m\ge1/16$;
the gain tomorrow in each state is $\sqrt A - \sqrt B \ge \tfrac{29}{10}d/(2\sqrt A) \ge \tfrac{29}{10}d/(2\cdot\tfrac{14}{25})$
since $A \le \tfrac{29}{160}+\tfrac18 \le (14/25)^2$; and $\tfrac{9}{10}\cdot\tfrac{29}{10}/\tfrac{28}{25} =
\tfrac{261}{112} > 2$. So Lemma~\ref{lem:hhi_of_le_policy} gives the upper condition, and
Theorem~\ref{thm:aiyagari_ge} applies; $K = L\,k(r) > 0$. Note $\beta(1+r_{hi}) > 1$, so the asset cap is what keeps
saving bounded at $r_{hi}$.
\end{proof}

\section{An endogenous asset bound}
\label{sec:asset_bound}

Throughout: a savings problem as in Definition~\ref{def:stoch_savings} with cash on hand $h(a,z) = Ra + y(z)$,
$R>1$, $0\le y(z)\le\bar y$, i.i.d.\ shocks (all rows of $Q$ equal), optimal policy $g$ (unique, continuous) and
consumption $c(a,z) = h(a,z) - g(a,z)$.

\begin{lemma}[Concave swap and antitone increments]
\label{lem:concave_swap}
\lean{Econ.ConcaveOn.add_le_add_of_between, Econ.ConcaveOn.sub_le_sub_of_le}
\leanok
If $f$ is concave, $p\le q_1\le r$ and $q_1+q_2 = p+r$, then $f(p)+f(r)\le f(q_1)+f(q_2)$. Hence for $p\le q$ and
$e\ge0$: $f(q)-f(q-e)\le f(p)-f(p-e)$.
\end{lemma}
\begin{proof}
\leanok
If $p<r$, write $q_1 = \lambda p + (1-\lambda) r$ and $q_2 = (1-\lambda)p + \lambda r$ and add the two concavity
inequalities. The second claim is the first with points $p-e\le p\le q$ and $q-e$.
\end{proof}

\begin{definition}[Euler gap]
\label{def:euler_gap}
\lean{Econ.EulerGap}
\leanok
$u$ has the Euler gap at $c$ with income gap $y$ (given $\beta$, $R$) if for every $\varepsilon>0$ there is
$d\in(0,\varepsilon]$ with $Rd\le c-y$ and $\beta[u(c-y) - u(c-y-Rd)] < u(c+d) - u(c)$: a difference form of
$\beta R\,u'(c-y) < u'(c)$.
\end{definition}

\begin{lemma}[Optimality tools]
\label{lem:asset_bound_tools}
\lean{Econ.StochSavings.cons, Econ.StochSavings.obj, Econ.StochSavings.feasible, Econ.StochSavings.cons_nonneg,
  Econ.StochSavings.obj_policy, Econ.StochSavings.obj_le, Econ.StochSavings.eq_of_swap,
  Econ.StochSavings.sum_eq_of_row, Econ.StochSavings.extendS_coe, Econ.StochSavings.value_concaveOn,
  Econ.StochSavings.value_sub_le, Econ.StochSavings.dev_le}
\leanok
\uses{def:stoch_savings, thm:stoch_value_concave}
(i) If alternative feasible choices at two states do at least as well in sum as the optimal ones, the first
alternative is optimal, hence equals the policy. (ii) Envelope: lowering cash on hand by $D\le c$ lowers the value by
at most $u(c)-u(c-D)$. (iii) Deviation: saving $y\le g(s)$ raises utility today by at most
$\beta\sum_{z'}Q(z,z')[V(g(s),z') - V(y,z')]$.
\end{lemma}
\begin{proof}
\leanok
(i) From the Bellman equation the alternative attains $V$, so it lies in the (singleton) argmax. (ii) From the
lower-cash state, keep the same saving and consume $D$ less. (iii) Compare the Bellman equation at $g(s)$ with
feasibility of $y$.
\end{proof}

\begin{theorem}[Monotone savings and consumption]
\label{thm:policy_cons_mono}
\lean{Econ.StochSavings.policy_mono, Econ.StochSavings.cons_mono}
\leanok
\uses{lem:asset_bound_tools, lem:concave_swap}
For two states with the same transition row and $h(s_1)\le h(s_2)$: $g(s_1)\le g(s_2)$ and $c(s_1)\le c(s_2)$.
\end{theorem}
\begin{proof}
\leanok
Savings: if $g(s_2)<g(s_1)$, swap the choices; concavity of $u$ (Lemma~\ref{lem:concave_swap}) makes the swap
weakly better, so by Lemma~\ref{lem:asset_bound_tools}(i) $g(s_2) = g(s_1)$. Consumption: if $c_2<c_1$, choose
$h_1 - c_2$ at $s_1$ and $h_2-c_1$ at $s_2$ (consumption swapped); utility today is unchanged and concavity of
$W(y) = \sum_{z'}Q(\cdot,z')V(y,z')$ makes the continuation weakly better, a contradiction as before.
\end{proof}

\begin{lemma}[Euler inequality and its two uses]
\label{lem:asset_bound_steps}
\lean{Econ.StochSavings.euler_ineq, Econ.StochSavings.cons_le_of_le_policy, Econ.StochSavings.cons_pos}
\leanok
\uses{lem:asset_bound_tools, thm:policy_cons_mono, def:euler_gap}
(a) If tomorrow's consumption is at least $p\ge Rd$ in every shock and $d\le g(s)$, then
$u(c+d) - u(c) \le \beta[u(p) - u(p-Rd)]$. (b) If $a\ge \bar y/R$, $a>0$, $g(a,z)\ge a$ and the Euler gap with income
gap $\bar y$ holds above $C^*$, then $c(a,z)\le C^*$. (c) Under the Inada condition
($u(d)-u(0) > Kd$ for small $d$, any $K$), $g(s)>0$ implies $c(s)>0$.
\end{lemma}
\begin{proof}
\leanok
(a) Save $d$ less today and consume $Rd$ less tomorrow: Lemma~\ref{lem:asset_bound_tools}(ii),(iii) and antitone
increments of $u$. (b) The state $(a-\bar y/R, z)$ has cash $h - \bar y$, so by monotone savings its consumption is
at least $c-\bar y$; tomorrow's states have cash at least $Ra\ge h-\bar y$, so by monotone consumption (i.i.d.\
shocks) consumption tomorrow is at least $c-\bar y$. If $c>C^*$, (a) with $p = c-\bar y$ contradicts the Euler gap.
(c) If $c=0$ and $g = y_1>0$, saving $d\le y_1/4$ less gains $u(d)-u(0)$ and loses at most $Kd$ by the chord-slope
bound for the concave $V(\cdot,z')$ on $[0,y_1/2]$ (\texttt{ConcaveOn.slope\_anti\_adjacent}).
\end{proof}

\begin{theorem}[Endogenous asset bound; Schechtman--Escudero, Aiyagari]
\label{thm:asset_bound}
\lean{Econ.StochSavings.policy_lt_self, Econ.StochSavings.policy_le_max, Econ.StochSavings.policy_lt_kbar,
  Econ.StochSavings.measure_gt_eq_zero}
\leanok
\uses{lem:asset_bound_steps, thm:policy_cons_mono, thm:stoch_savings_stationary}
Assume the Inada condition, the Euler gap with gap $0$ for all $c>0$ and with gap $\bar y$ for $c>C^*$. If
$A>0$, $A\ge \bar y/R$ and $(R-1)A\ge C^*$, then $g(a,z)<a$ whenever $a\ge A$; $g(a,z)\le\max(A,a)$; if $A<\bar k$
then $g<\bar k$ everywhere (the cap is slack); and every stationary distribution of the policy kernel gives
$\{a>A\}$ measure zero. The bound $A$ does not depend on $\bar k$.
\end{theorem}
\begin{proof}
\leanok
Suppose the closed set $N = \{(a,z): a\ge A,\ g(a,z)\ge a\}$ is nonempty and pick $s_0\in N$ minimizing $c$. By
Lemma~\ref{lem:asset_bound_steps}(b),(c), $0<c_0\le C^*$. Tomorrow, from $a_1 = g(s_0)\ge A$: if $(a_1,z')\in N$ its
consumption is $\ge c_0$ by minimality; otherwise $g(a_1,z')<a_1$ and consumption exceeds $(R-1)a_1\ge C^*\ge c_0$.
Lemma~\ref{lem:asset_bound_steps}(a) with $p = c_0$ contradicts the Euler gap at gap $0$. Below $A$, monotone
savings give $g(a,z)\le g(A,z)<A$. For stationary $\pi$, $\varphi(a) = \max(a-A,0)$ satisfies
$\int\varphi\,d\pi = \int\varphi\circ g\,d\pi$ while $\varphi(g)\le\varphi(a)$, strictly when $a>A$.
\end{proof}

\begin{lemma}[Square-root utility]
\label{lem:sqrt_euler_gap}
\lean{Econ.sqrt_inada, Econ.sqrt_eulerGap}
\leanok
\uses{def:euler_gap}
$u=\sqrt{\cdot}$ satisfies the Inada condition, and if $0<\beta R<1$ and $y<(1-(\beta R)^2)c$ it has the Euler gap at
$c$ with income gap $y$. So $C^* = \bar y/(1-(\beta R)^2)$ works, and gap $0$ holds for every $c>0$.
\end{lemma}
\begin{proof}
\leanok
Inada: $d = t^2$ with $t$ small. Euler gap: at $d=0$, $\sqrt{c-y}+\sqrt{c-y} > \beta R(\sqrt c+\sqrt c)$ since
$c-y>(\beta R)^2c$; by continuity this persists for small $d>0$, and with
$\sqrt{c+d}-\sqrt c = d/(\sqrt{c+d}+\sqrt c)$ and the analogous identity it gives the gap.
\end{proof}

\begin{theorem}[Examples with income risk where the cap is slack]
\label{thm:asset_bound_examples}
\lean{Econ.StochSavings.riskyHousehold, Econ.StochSavings.riskyHousehold_assetBound,
  Econ.AiyagariEconomy.household_cap_slack, Econ.AiyagariEconomy.riskyEconomy,
  Econ.AiyagariEconomy.risky_cap_slack}
\leanok
\uses{thm:asset_bound, lem:sqrt_euler_gap, def:aiyagari_economy}
(1) $u=\sqrt c$, $R = 11/10$, $\beta=1/2$, i.i.d.\ income $y\in\{0,1\}$, $\bar k=20$: with $C^*=3/2$, $A = 15$, assets
above $15$ fall, $g<20$, a stationary distribution exists and every stationary distribution lives on $[0,15]$. (2) For an Aiyagari economy at $r>0$ the same
holds with $\bar y = w(r)\bar e$ and $(1+r)$ in place of $R$. In the economy with $e\in\{0,2\}$ i.i.d., $\beta=1/2$,
$\alpha=1/2$, $\delta=1/10$, $\bar k = 50$, at $r = 1/10$: $w = 5/4$, $\bar y = 5/2$, $C^* = 4$, $A = 40<50$.
\end{theorem}
\begin{proof}
\leanok
Instantiate Theorem~\ref{thm:asset_bound} with Lemma~\ref{lem:sqrt_euler_gap}: $1 < \tfrac{279}{400}\cdot\tfrac32$ and
$\tfrac52<\tfrac{279}{400}\cdot4$; $(R-1)A = C^*$ in both cases.
\end{proof}

\noindent\emph{Not covered.} $r\le0$; persistent (Markov) income; and a general-equilibrium witness in which both
boundary conditions of Theorem~\ref{thm:aiyagari_ge} are derived \emph{and} the cap is slack (this needs the
divergence of supply as $\beta(1+r)\uparrow1$).
